% Imporditud: tools/import_eksperimendid.py
% Allikas: Eesti füüsikaolümpiaadi eksperimendi ülesannete
%   kogu (Rael Kalda, 2023), ülesanne Ü163, lahendus L163.
\setAuthor{EFO žürii}
\setRound{lõppvoor}
\setYear{2003}
\setNumber{G E2}
\setDifficulty{5}
\setTopic{elektriahelad}
\setVahendid{tester (multimeeter), lapik patarei, juhtmed, käärid, kuubikujuline plastanum, tükk paksu alumiiniumfooliumi, vesi, mõõtejoonlaud, paberkäterätik töökoha kuivatamiseks}
\setPunktid{14}
\setAllikas{Eksperimendi kogu Ü163, lahendus lk 122}
\setLahendusseis{puudub}

\prob{Vee eritakistus}
Määrake kraanivee eritakistus. Visandage protokolli katseskeem(id), millelt on näha elektroodide kuju ja paigutus vees ning kõik ühendused.
\textit{Märkus:} ed: Mitte kasutada testrit oommeetrina. Arvestada, et kasutatava testri takistus voltmeetrina mõõtepiirkonnal 20 V on 1 M$\Omega$, ampermeetrina mõõtepiirkonnal 200 mA --- 6 $\Omega$, piirkonnal 20 mA --- 15 $\Omega$ ja piirkonnal 2 mA --- 105 $\Omega$. Eeldada, et patarei sisetakistus ei ületa ühte oomi. Testri plussklemm on parempoolne pistikupesa, miinusklemm aga keskmine pesa. Töölõpetamisel lülitada tester kindlasti välja (asend OFF).

\hint

\solu
\textit{Kogumikus lahendus puudub.}
\probend
